% 流网络与割：容量 s->a 3, s->b 2, a->b 1, a->t 2, b->t 3；最大流各弧 3,2,1,2,3，流值 5。
% 割 S={s,b}：跨出 s->a(3)+b->t(3)=6；指回 a->b 流量 1；净流量 3+3-1=5。
\begin{tikzpicture}[>=Stealth,
  v/.style={circle,draw=FigureStroke,fill=white,line width=.8pt,minimum size=6.5mm,inner sep=0pt},
  arc/.style={->,draw=FigureStroke,line width=.8pt},
  fwd/.style={->,draw=FigureBlue,line width=1.6pt},
  back/.style={->,draw=FigureRed,dashed,line width=1.2pt},
  w/.style={fill=white,inner sep=.8pt,font=\figurefont\fontsize{7.5}{9}\selectfont},
  note/.style={font=\FigureNoteFont\fontsize{7.5}{9}\selectfont}]
  \fill[FigureBlueFill,rounded corners=4mm] (-.55,.55)--(-.55,-.5)--(1.6,-1.65)--(2.6,-1.65)--(2.6,-.55)--(.5,.55)--cycle;
  \node[v] (s) at (0,0) {$s$};
  \node[v] (a) at (2,1.1) {$a$};
  \node[v] (b) at (2,-1.1) {$b$};
  \node[v] (t) at (4,0) {$t$};
  \draw[fwd] (s)--node[w]{$3/3$}(a);
  \draw[arc] (s)--node[w]{$2/2$}(b);
  \draw[back] (a)--node[w]{$1/1$}(b);
  \draw[arc] (a)--node[w]{$2/2$}(t);
  \draw[fwd] (b)--node[w]{$3/3$}(t);
  \node[note,text=FigureBlue] at (.75,-1.45) {$S=\{s,b\}$};
  \node[note,align=left,anchor=west] at (4.5,.45) {割容量$c(S,\bar S)=3+3=6$};
  \node[note,align=left,anchor=west] at (4.5,-.05) {净流量$3+3-1=5$};
  \node[note,align=left,anchor=west] at (4.5,-.55) {流值$\operatorname{val}(\bm x)=5$};
\end{tikzpicture}
